% =====================================================================
\section{Introduction}
\label{sec:intro}
% =====================================================================

Large language models (LLMs) are increasingly composed into multi-stage
pipelines in which each stage is an agent with a typed output contract, and
those pipelines are shipped as software components. Structured-output tooling
has made the components reliably \emph{well-formed}, but well-formed is not
grounded. An evidence identifier can be syntactically perfect and point at
nothing. A recommendation can validate against its schema and have no
connection to anything an earlier stage found. A report can restate content
that no stage produced. Each of these passes every check the type system
offers. Catching such defects while the pipeline runs, and doing something
about them, is a verification problem in the ordinary software-engineering
sense: a component's output must be checked against a specification that its
types do not capture~\citep{leucker2009rv}.

In the conference version of this work~\citep{hoang2026eas}, we built such a
check into a four-stage pipeline that analyses restaurant reviews. Every stage
reads only what earlier stages produced in the same run, so every claim a stage
makes can be checked against upstream state alone. A deterministic supervisor
measures each stage's output, decides by fixed rules whether to proceed, warn,
retry or halt, and on a retry writes the measured failure into the prompt of
the stage being re-run. It makes no model call of its own. That study measured
how often an unsupervised pipeline emits ungrounded content. It could not
measure whether the retry \emph{repairs} it: its enforcing arm produced a single
retry. Repair is the property that determines whether such a supervisor is
worth building, because a runtime check that detects defects but cannot get
them fixed only halts the pipeline.

Whether measured feedback repairs is genuinely contested, and the two relevant
literatures disagree in a way an experiment can settle. One holds that feedback
\emph{content} is the active ingredient. In code repair, better feedback
produces better repairs. On 40 failing programs, human-written explanations
raised the share of successful repairs from 33.3\% to 52.6\% over model-written
ones~\citep{olausson2024selfrepair}. Models that cannot find their own reasoning
errors can correct many of them once the location is
supplied~\citep{tyen2024}. The other holds
that the \emph{signal} that something is wrong does most of the work. A bare
continuation cue closes most of the gap between how models correct others'
errors and their own~\citep{tsui2026scbench}. In self-debugging, the least
informative feedback condition recovers much of the gain of the most
informative one when execution still decides whether a revision is
needed~\citep{chen2024selfdebug}. A deterministic supervisor's feedback sits at
the informative extreme. It names the offending item and the measured fact.
If content binds, measured-fact feedback should repair more defects than a
rejection that says nothing. If the signal is what matters, the two should
repair alike.

We test this in a pre-registered fault-injection study. Six operators corrupt
clean stage outputs taken from 28 real pipeline runs, giving 142 defects, each
designed to trip exactly one supervisor check. Each defective stage is re-run
through its production agent under four arms: the production measured-fact
feedback; the same message with the facts removed; the defect presented as a
fresh correction task, which serves as a capability ceiling; and a bare
regeneration. The single confirmatory test, fixed with the arm wording, the
outcome coding and the analysis before any trial ran, compares the first two.
Around the confirmatory core of 1{,}167 trials, exploratory cells vary
where the defective output appears, the sampling temperature, the extent of
the defect, whether the stage's obligations are restated, and the reasoning
budget allocated to internal thinking traces. Table~\ref{tab:cells} distinguishes trials from downstream conditions and calls. They also replay 13 real defects that the production
supervisor caught, follow every defect into the later stages, and measure what
the supervisor's deterministic overwrite of one statistic buys over a retry.
Another repeats the core's arms on a second, open-weight model served locally.

The confirmatory contrast found little difference in overall repair on the
hosted model near its ceiling. Exploratory analyses found that measured facts
increased unconditional exact restoration to the clean output, and that
restoration outcomes were associated with downstream divergence. These
post-treatment associations do not establish a causal mechanism. The local
model showed lower repair and an exploratory feedback advantage. The production
retry was accepted much less often on real defects than on injected ones, many
of which were judged differently by blinded raters.

The contributions are a pre-registered, gated fault-injection protocol and
paired feedback-content comparison; evidence on exact restoration and its
qualified downstream association; a real-defect replay and second-model
comparison; and calibration and propagation analyses that expose limits of
stage-local verification. Section~\ref{sec:design-rq} states RQ1--RQ6.

\paragraph{Relation to the conference version}
\label{sec:intro-extension}
This article extends the conference paper accepted at EAI FISAT
2026~\citep{hoang2026eas}, which introduced the supervisor, the derivability
property and the measurements of defect incidence and instability.
Section~\ref{sec:system} summarises that system in new text. Everything from
Section~\ref{sec:design} onwards is new: the pre-registered repair study, every
trial, the real-defect replay, the downstream-propagation analysis and the
calibration. Section~\ref{sec:discussion-reconcile} lists the points at which
the new results revise claims made in the conference version.

The rest of the paper is organised as follows. Section~\ref{sec:related}
positions the study. Section~\ref{sec:system} describes the system and
Section~\ref{sec:design} the study design. Section~\ref{sec:results} reports
the results by research question. Section~\ref{sec:discussion} discusses what
they mean for building runtime supervisors, Section~\ref{sec:threats} the
threats to validity, and Section~\ref{sec:conclusion} concludes.
